\chapter{Grafovska predstavitev podatkov sledilnika pogleda}
\label{sec:grafovska-predstavitev}

V prejšnjem poglavju smo meritve sledilnika pogleda razdelili na 10-sekundna časovna okna.
V tem poglavju opišemo, kako posamezno okno pretvorimo v graf, ki ga uporabimo kot vhod za grafovske modele.
Vozlišča predstavljajo veljavne meritve znotraj okna, povezave pa izbrane odnose med meritvami.

Najprej podamo najbogatejšo grafovsko predstavitev učnega vzorca, ki uporablja vse razpoložljive signale.
Nato pokažemo, kako se ista osnovna konstrukcija prilagodi štirim množicam signalov iz poglavja~\ref{subsec:vhodni-signali}.
Kadar določena množica signalov ne vsebuje informacije, potrebne za posamezno značilko ali relacijo, ta del predstavitve izpustimo.
Graf zato v vseh primerih predstavlja isto vrsto učnega vzorca, razlikuje pa se po količini vhodne informacije in po razpoložljivih relacijah med meritvami.

\section{Vozlišča}
\subsection{Definicija}

Vsako vozlišče predstavlja eno meritev sledilnika pogleda znotraj časovnega okna:
\begin{equation}
    v_i \leftrightarrow \mathbf{x}_i.
\end{equation}
Če okno vsebuje $n$ veljavnih meritev, potem graf vsebuje $n$ vozlišč.
Vrednost $n$ se lahko med okni razlikuje, saj po čiščenju in filtriranju niso vsa okna popolna.
Grafovska predstavitev naravno podpira spremenljivo število vozlišč.

\subsection{Vektor značilk vozlišča}
Vektor značilk vozlišča za končni graf za množico \emph{vsi signali} je
\begin{equation}
    \mathbf{x}_i =
    \left[
        t_i,
        x_i,
        y_i,
        p_i^{(l)},
        p_i^{(r)},
        d_i,
        f_{dur,i}
    \right],
\end{equation} % fact check the equation
kjer je $t_i$ normaliziran časovni žig za vozlišče $i$, $x_i$ in $y_i$ sta koordinati pogleda, $p_i^{(l)}$ in $p_i^{(r)}$ velikosti leve in desne zenice, $d_i$ je razdalja med sledilnikom in udeleženčevimi očmi, $f_{dur,i}$ pa trajanje fiksacije.

Za ostale podgrafe se vektor značilk ustrezno filtrira.
V vseh primerih ohranimo časovno značilko $t_i$, ostale podatke pa ohranimo glede na obravnavano množico signalov.

\section{Relacije}
Definiramo tri relacije med meritvami: časovno, prostorsko in fiksacijsko.
Posledično v grafu oblikujemo štiri tipe povezav.
Pri tem sta časovna tipa povezav usmerjena -- časovne povezave naprej in nazaj.
Prostorske in fiksacijske povezave pa obravnavamo kot neusmerjene.

\subsection{Zakaj ločimo relacije}

Vsaka povezava ima določen tip relacije, saj časovna povezava, prostorska povezava in fiksacijska povezava opisujejo različne odnose med meritvami.
Ta zapis omogoča, da kasnejše stopnje arhitekturne lestvice relacije obravnavajo ločeno.
Osnovni GCN uporabi isti nabor povezav, vendar relacij modelno ne loči; v tem primeru dobimo enostaven homogen graf.
Podrobnosti o uporabi relacij v posameznih modelih so opisane v poglavju~\ref{sec:modeli}.

\subsection{Časovne povezave}
Časovne povezave povezujejo meritve zaporedne v času.
Naj bodo vozlišča v posameznem oknu urejena po času, tako da je $v_i$ $i$-ta meritev v časovnem zaporedju.

Ločimo usmerjene časovne povezave naprej in nazaj:
\begin{align}
    E_t^+ &= \{(v_i, v_j) : 1 \leq j - i \leq k_t\}, \\
    E_t^- &= \{(v_i, v_j) : 1 \leq i - j \leq k_t\}.
\end{align}

Parameter $k_t \in \mathbb{N}$ določa število časovnih sosedov na vsako stran.
Trenutna konfiguracija uporablja $k_t=1$, zato je vozlišče $v_i$ povezano samo z neposredno predhodno in neposredno naslednjo meritvijo, kadar ti vozlišči obstajata.

Povezave v $E_t^+$ potekajo od prejšnjih proti kasnejšim meritvam, povezave v $E_t^-$ pa od kasnejših proti prejšnjim meritvam.
Pri grafu z $n$ vozlišči to pri $k_t=1$ pomeni $n-1$ časovnih povezav naprej in $n-1$ časovnih povezav nazaj.
Za polno 10-sekundno okno s pričakovanimi $600$ meritvami je to $599$ povezav v vsako smer.


\subsection{Prostorske povezave}
\label{subsec:prostorske-povezave}

Prostorske povezave so neusmerjene povezave, ki povezujejo meritve s podobnim položajem pogleda znotraj istega časovnega okna.
Te povezave lahko definiramo samo za grafe, ki imajo prostorsko informacijo, torej signala $(x,~y)$
Za vsako vozlišče poiščemo $k_s$ najbližjih sosedov v prostoru koordinat pogleda $(x,~y)$.

Prostorsko razdaljo med meritvama $i$ in $j$ definiramo evklidsko:
\begin{equation}
    d_{ij}^{(s)} =
    \sqrt{(x_i - x_j)^2 + (y_i - y_j)^2}.
\end{equation}

Za vsako vozlišče $v_i$ poiščemo $k_s$ najbližjih sosedov v prostoru koordinat pogleda $(x,~y)$ z uporabo k-d drevesa. Če je $v_j$ med temi sosedi, dodamo povezavo med vozliščema $v_i$ in $v_j$.
Ker prostorske povezave obravnavamo kot neusmerjene, povezavo predstavimo z dvema usmerjenima povezavama $(v_i,v_j)$ in $(v_j,v_i)$.
Na koncu odstranimo podvojene povezave znotraj relacije $E_s$.

Število prostorskih povezav zato ni določeno samo z $n$ in $k_s$, ampak tudi z geometrijo položajev pogleda.
Pred simetrizacijo vsako vozlišče prispeva $k_s$ sosedskih izbir, po simetrizaciji in odstranitvi podvojenih parov pa za usmerjeni zapis neusmerjene relacije velja
\begin{equation}
    n k_s \leq |E_s| \leq 2 n k_s.
\end{equation}
Dejansko gostoto prostorskih povezav zato poročamo z opazovanimi statistikami grafov.

\subsection{Fiksacijske povezave}
\label{subsec:fiksacijske-povezave}

Fiksacijske povezave povezujejo meritve, ki pripadajo isti zaznani fiksaciji.
Pripadnost fiksaciji določimo z identifikatorjem fiksacije, podanim v \MAHNOB podatkih.
Ta podatek imamo zgolj v množici \emph{vsi signali}, zato so te povezave vključene le v končni graf.

\subsubsection{Redčenje povezav znotraj fiksacije}

Znotraj posamezne fiksacije ne povežemo vseh vozlišč med seboj, saj bi to povzročilo preveliko število fiksacijskih povezav.
Namesto tega uporabimo redkejšo konstrukcijo, ki ohrani informacijo o pripadnosti isti fiksaciji, hkrati pa zmanjša prostorsko in časovno zahtevnost konstrukcije grafa ter učenja na grafu.
To konstrukcijo sami imenujemo \emph{razširjene povezave znotraj fiksacije} (angl. \emph{dilated intra-fixation edges}).
Fiksacijo obravnavamo kot neprekinjeno zaporedje vozlišč z istim veljavnim identifikatorjem \texttt{fixation-index}.
Če ima taka fiksacija $F$ vozlišč, jim znotraj fiksacije priredimo lokalne indekse $0,\ldots,F-1$.

Gostoto povezav uravnava parameter $k_f$, ki določa število fiksacijskih povezav, ki jih generira vsako vozlišče.
V trenutni konfiguraciji uporabljamo $k_f=3$.
Za posamezno fiksacijo najprej izračunamo korak razširitve
\begin{equation}
    s = \max(1, \lfloor F/k_f + 0{,}5 \rfloor).
\end{equation}
Množico kandidatnih odmikov nato definiramo kot
\begin{equation}
    D = \{(1 + q s) \bmod F \mid q = 0,\ldots,k_f-1\}.
\end{equation}
Morebitni odmik $0$ in podvojene odmike odstranimo.
Za vsak odmik $d \in D$ dodamo neusmerjene povezave
\begin{equation}
    \{v_i, v_{(i+d)\bmod F}\},
    \qquad i = 0,\ldots,F-1.
\end{equation}
Pri zapisu grafa za posredovanje sporočil vsako neusmerjeno povezavo predstavimo z dvema usmerjenima povezavama.
Podvojene povezave znotraj relacije $E_f$ odstranimo.

Na primer, pri fiksaciji dolžine $F=30$ in parametru $k_f=3$ dobimo $s=10$.
Vozlišče z lokalnim indeksom $0$ torej povežemo z vozlišči $1, 11$ in $21$.
Odmik $1$ ohrani časovno lokalno komunikacijo znotraj fiksacije, večji odmiki pa omogočijo širjenje informacij tudi med časovno bolj oddaljenimi meritvami iste fiksacije.

Razširjene povezave imajo večji doseg kot zgolj zaporedne fiksacijske povezave, hkrati pa se izognejo kombinatorični eksploziji, ki bi nastala pri polnem kliku.
Pri polnem kliku bi namreč število povezav naraščalo kvadratno z dolžino fiksacije.
Pri analiziranih 10-sekundnih oknih dobimo približno $1.956$ fiksacijskih povezav na graf pri $k_f=2$, $2.908$ pri $k_f=3$, $4.163$ pri $k_f=5$ in $7.273$ pri $k_f=10$.
Pri privzeti izbiri $k_f=3$ so fiksacijske povezave predstavljale približno $45{,}7\,\%$ vseh povezav v grafu (relacije \texttt{temporal}, \texttt{spatial} in \texttt{fixation}).
%TODO: fact check številke

Izbira $k_f=3$ je zato kompromis med dosegom in gostoto grafa.
Ker vsaka plast GNN omogoča izmenjavo informacij med neposredno povezanimi vozlišči, se doseg znotraj fiksacije povečuje tudi z večanjem števila plasti. % omeni kako se povečuje (če je to enostavno reči) -- kvadratno, eksponentno, ...?
Pri $k_f=3$ vozlišče že v eni plasti prejme informacije iz več delov iste fiksacije, večplastni model pa lahko te informacije dodatno propagira po fiksacijski relaciji.
Večje vrednosti $k_f$ povečajo gostoto grafa, računsko zahtevnost in možnosti za prekomerno glajenje, manjše vrednosti pa graf dodatno razredčijo.

\subsection{Značilke povezav}
\label{subsec:utezi-povezav}

Povezavam poleg tipa relacije priredimo tudi značilke povezav, ki opisujejo odnos med izvorno meritvijo $v_i$ in ciljno meritvijo $v_j$.
Osnovne časovne značilke so normalizirani čas izvorne in ciljne meritve, $t_i$ in $t_j$, ter časovna razlika $\Delta t = t_j - t_i$.
Pri časovnih povezavah dodamo še smer $\rho_{ij}$, kjer ima povezava naprej vrednost $\rho_{ij}=1$, povezava nazaj pa $\rho_{ij}=-1$.
Značilke povezav so del najbogatejše grafovske predstavitve; eksplicitno jih uporabi utežena stopnja arhitekturne lestvice, opisana v poglavju~\ref{sec:modeli}.

Pri grafih, ki vsebujejo koordinate pogleda, povezavam dodamo tudi prostorske značilke.
To sta razliki v koordinatah pogleda, $\Delta x = x_j - x_i$ in $\Delta y = y_j - y_i$, ter evklidska razdalja med položajema pogleda na zaslonu $d_{ij}^{screen}$.
Te značilke so zato na voljo pri vseh množicah signalov, razen pri množici \emph{samo zenici}.

Pri množici \emph{vsi signali} povezavam dodamo še značilko, povezano z oddaljenostjo od sledilnika pogleda.
Značilka $\Delta d^{eye}_{ij}$ opisuje razliko v oddaljenosti med ciljno in izvorno meritvijo.

Pregled vseh značilk povezav in njihove uporabe je podan v tabeli~\ref{tab:znacilke-povezav}.

\begin{table}[H]
    \centering
    \caption{Značilke povezav v grafovski predstavitvi.}
    \label{tab:znacilke-povezav}
    \begin{tabularx}{\textwidth}{@{}>{\raggedright\arraybackslash}p{0.18\textwidth} >{\raggedright\arraybackslash}X >{\raggedright\arraybackslash}p{0.30\textwidth}@{}}
        \toprule
        Značilka & Opis & Uporaba \\
        \midrule
        $t_i$ & Normaliziran čas izvorne meritve znotraj časovnega okna. & Vse povezave. \\
        $t_j$ & Normaliziran čas ciljne meritve znotraj časovnega okna. & Vse povezave. \\
        $\Delta t$ & Časovna razlika med ciljno in izvorno meritvijo, $\Delta t = t_j - t_i$. & Vse povezave. \\
        $\rho_{ij}$ & Smer časovne povezave; pri povezavi naprej je $\rho_{ij}=1$, pri povezavi nazaj pa $\rho_{ij}=-1$. & Samo časovne povezave. \\
        $\Delta x$ & Razlika v vodoravni koordinati pogleda, $\Delta x = x_j - x_i$. & Grafi s koordinatami pogleda. \\
        $\Delta y$ & Razlika v navpični koordinati pogleda, $\Delta y = y_j - y_i$. & Grafi s koordinatami pogleda. \\
        $d_{ij}^{screen}$ & Evklidska razdalja med položajema pogleda na zaslonu. & Grafi s koordinatami pogleda. \\
        $\Delta d^{eye}_{ij}$ & Razlika v oddaljenosti od sledilnika pogleda med ciljno in izvorno meritvijo. & Samo \emph{vsi signali}. \\
        \bottomrule
    \end{tabularx}
\end{table}


\section{Grafovska konstrukcija po množicah signalov}
\label{subsec:grafovska-konstrukcija-po-mnozicah}

Za vsako množico signalov iz tabele~\ref{tab:mnozice-signalov} zgradimo graf po isti osnovni konstrukciji.
Razlika med grafi je v značilkah vozlišč in relacijah, ki jih izbrana množica signalov omogoča.
Množica \emph{vsi signali} je najbogatejša predstavitev. 
Ostale množice uporabijo ustrezne podmnožice značilk in povezav.

\begin{table}[H]
    \centering
    \caption{Grafovska konstrukcija po množicah signalov.}
    \label{tab:grafovska-konstrukcija-po-mnozicah}
    \begin{tabularx}{\textwidth}{@{}>{\raggedright\arraybackslash}p{0.20\textwidth} >{\raggedright\arraybackslash}X >{\raggedright\arraybackslash}X@{}}
        \toprule
        Množica signalov & Značilke vozlišč & Tipi povezav \\
        \midrule
        \emph{samo pogled} & čas in koordinati pogleda & časovne naprej, časovne nazaj, prostorske \\
        \emph{samo zenici} & čas in velikosti obeh zenic & časovne naprej, časovne nazaj \\
        \emph{pogled in zenici} & čas, koordinati pogleda in velikosti obeh zenic & časovne naprej, časovne nazaj, prostorske \\
        \emph{vsi signali} & čas, koordinati pogleda, velikosti obeh zenic, oddaljenost od sledilnika pogleda in fiksacijska informacija & časovne naprej, časovne nazaj, prostorske, fiksacijske \\
        \bottomrule
    \end{tabularx}
\end{table}

Slika~\ref{fig:konstrukcija-grafa-primer} prikazuje primer najbogatejše konstrukcije na majhnem izseku časovnega okna za množico \emph{vsi signali}.
Prikaz združi vozlišča, časovne, prostorske in fiksacijske povezave v isti graf.

\begin{figure}[H]
    \centering
    \includegraphics[width=\textwidth]{gnn_graph_far_02_P10_90_avi_t106_86-107_95_kt1_ks1_kf3_f3_n9/vse_povezave.png}
    \caption{Primer grafovske konstrukcije za množico \emph{vsi signali}. Vozlišča predstavljajo meritve v izseku časovnega okna, povezave pa prikazujejo časovne, prostorske in fiksacijske relacije.}
    \label{fig:konstrukcija-grafa-primer}
\end{figure}

Ločeni prikazi posameznih relacij so zbrani v dodatku~\ref{app:grafi}.

\section{Statistika konstruiranih grafov}
\label{subsec:statistika-grafov}

% Namen sekcije:
% - Zaključiti poglavje z lastnostmi grafov, ki dejansko nastanejo iz podatkov.
% - To ni razdelek o zahtevnosti modelov.
% Vsebina:
% - Preračunaj statistiko za končno konfiguracijo in za vse štiri množice signalov.
% - Tabela naj verjetno vsebuje stolpce:
%   - \emph{samo pogled};
%   - \emph{samo zenici};
%   - \emph{pogled in zenici};
%   - \emph{vsi signali}.
% - Vrstice naj bodo na primer:
%   - mediana števila vozlišč;
%   - povprečno število vozlišč;
%   - mediana časovnih povezav naprej;
%   - mediana časovnih povezav nazaj;
%   - mediana prostorskih povezav;
%   - mediana fiksacijskih povezav;
%   - mediana skupnega števila povezav.
% - Pri relacijah, ki v določeni množici signalov ne obstajajo, uporabi nič; pojasni v opombi pod tabelo.
% - Preveri skladnost s končnimi parametri `kt`=1, `ks`=1, `kf`=3.
% Uporabi iz stare verzije:
% - Staro tabelo kot oblikovni začetek, ne kot vir končnih številk.
% Ne uporabi iz stare verzije:
% - Odseka "Računska zahtevnost".
% - Primerjave števila parametrov, časa učenja, časa inference ali modelne zahtevnosti; to sodi v poglavje 7, 8 ali dodatek s hiperparametri.
